Sparse regression on a library of candidate terms is a popular route to recovering governing equations from data. SINDy~\cite{brunton2015discovering} regresses the time derivative of the state on a library and promotes sparsity with sequentially thresholded least squares; PDE-FIND~\cite{rudy2016data} extended the idea to partial differential equations. The derivative is not measured, so it has to be estimated from noisy samples, and the quality of that estimate often decides whether the correct terms are found. Toolboxes such as PySINDy~\cite{silva2020pysindy} therefore expose several differentiators, and the weak or integral formulation~\cite{schaeffer2017sparse,messenger2021weak,messenger2020weak,reinbold2020using} sidesteps differentiation altogether by integrating the dynamics against test functions.

The ODE Weak SINDy paper~\cite{messenger2021weak} compares the weak form with standard SINDy. This study asks a narrower, controlled question about all of these options: with the same library, the same STLSQ solver, the same data and the same tuning protocol, how do finite differences, Savitzky--Golay~\cite{savitzky1964smoothing}, a smoothing spline, TV-regularised differentiation~\cite{chartrand2011numerical} and the weak form compare in the \emph{rate of exact support recovery} and in coefficient error on Lorenz-63~\cite{lorenz1963deterministic} as the noise grows from 0 to 10\%?

We registered five hypotheses before the runs (Sec.~\ref{sec:setup}). The outcome is mostly negative: H1 (the weak form beats every derivative-based method in support rate from 1\% noise) and H2 (every smoother has a lower coefficient error than finite differences) are refuted by the registered tests, H2 with a verdict that depends on the finite-difference threshold at 5\%; H3 (all systems are exact at zero noise) is supported; H4 (too narrow a test-function window loses the advantage) is partly supported, since the width matters but in the opposite direction; H5 (a library on the smoothed state improves support recovery) is not supported. Contributions are (i) a controlled, reproducible comparison with failures split into false positives and negatives; (ii) a grid of every system at every STLSQ threshold, which shows that the low-noise ordering of the methods is decided by threshold selection, including the tie rule of our own tuner, revised after a first round of results; (iii) width, exponent and library-state ablations. Scale is small: one simulated system, 20 trajectories per cell, reimplemented methods, CPU.
